\documentclass[nonacm,sigconf]{acmart}
\usepackage{paralist}
\usepackage{xurl}
\setlength{\emergencystretch}{3em}
\Urlmuskip=0mu\relax
\settopmatter{printacmref=false,printfolios=true}
\renewcommand\footnotetextcopyrightpermission[1]{}
\pagestyle{plain}
\hypersetup{pdftitle={PhantomChannel Artifact Appendix},pdfauthor={Xu Zhang et al.}}
\DeclareUrlCommand\file{\urlstyle{tt}}
\begin{document}
\title{PhantomChannel: Artifact Appendix}
\author{Xu Zhang, Ming Zhou, Meng Zhang, Dong Wei, Haining Wang, and Qiaoyu Zhang}
\renewcommand{\shortauthors}{Xu Zhang et al.}
\maketitle
\appendix
\section{Artifact Appendix}

\subsection{Abstract}
PhantomChannel embeds auxiliary symbols beyond the protocol-visible boundary
of BLE packets and recovers them with an SDR receiver. The artifact contains
transmitter and receiver source, tests, 14 firmware images for five platform
families, two demonstration videos, and a self-contained offline experiment
bundle. The bundle regenerates plots and numerical summaries from analysis-ready
records. This appendix specifies the reproducible subset of the evaluation and
the metric definitions used to interpret its outputs. Project-owned source is
MIT licensed; videos are CC BY 4.0. Third-party code and firmware retain their
file-level and vendor terms (\file{THIRD_PARTY_NOTICES.md}).

\subsection{Description \& Requirements}
The source archive contains \file{receiver/}, \file{firmware/},
\file{patches/zephyr/}, \file{configs/}, \file{tools/}, and \file{tests/}.
The experiment archive contains \file{data/}, \file{scripts/}, \file{checks/},
\file{reproduce_all.py}, and a reverse dependency map in
\file{DEPENDENCY_MAP.md}. Generated CSV/JSON results are under
\file{outputs/}; plots appear in the root and \file{plots/}. Each archive
has its own README. The workflow reanalyzes retained records; it does not
repeat RF acquisition or independently reconstruct all upstream measurements.

\subsubsection{Security, privacy, and ethical concerns}
Offline analysis requires no privileged access or radio transmission. Hardware
evaluation flashes development boards and emits BLE traffic; flashing replaces
the existing firmware. Use owned devices and authorized, isolated RF setups
under applicable spectrum rules. Keep credentials and device identifiers out of
shared configurations and captures. Published packet records use anonymous
groups and relative timing.

\subsubsection{How to access}
The permanent artifact is available at
\href{https://doi.org/10.5281/zenodo.22842761}{doi:10.5281/\allowbreak zenodo.22842761}.
The record version is \emph{PhantomChannel V1.0}; the archive suffix is
\texttt{2026.09.19-ae}. Download the source and experiment ZIPs with that
suffix, the top-level README, and \file{SHA256SUMS}. Verify downloaded
files with \texttt{sha256sum --ignore-missing -c SHA256SUMS}.
The maintained repository is
\url{https://github.com/Tincyu/PhantomChannel}; the archived snapshot fixes
the evaluation version.

\subsubsection{Hardware dependencies}
Offline execution needs a conventional CPU workstation; no SDR, GPU, or BLE
board is required. Allow approximately 2\,GB of working disk space and 4\,GB RAM
as a practical budget. Optional live evaluation uses an nRF52840 DK transmitter,
a compatible BLE peer, and a USRP B210 or X310 receiver; the 80-MHz bandwidth
study uses X310. Alternative transmitters cover nRF52833, DA14695, EFR32MG24,
CC2340R5, and STM32L476RG with SX1280. Board/image requirements are in
\file{firmware/images/FIRMWARE_GUIDE.md}.

\subsubsection{Software dependencies}
Use Linux, Python 3.10+, \texttt{venv}, and \texttt{pip}. Source tests use
Python 3.12 with NumPy 2.4.6, SciPy 1.17.1, PyYAML 6.0.3, and pytest 8.4.2
as pinned in \file{requirements-test.txt}. Analysis declares NumPy
$\geq$1.26, SciPy $\geq$1.11, and Matplotlib $\geq$3.8; an audited run used
2.5.4, 1.18.1, and 3.11.2, respectively, on Python 3.12.3 / Ubuntu 24.04.4
LTS. Record the actual Python/OS and
\texttt{pip freeze} output for each run. Live reception additionally requires
UHD and the selected receiver runtime; board rebuilding requires external
vendor SDKs and toolchains. For Nordic, the guide specifies NCS 3.0.0-rc1
and the controller patch base \texttt{v4.0.99-ncs1-1}.

\subsubsection{Benchmarks}
All offline inputs are bundled: 14 complex IQ packet windows, 23,789 parser
candidate rows from 52 range runs, advertising timing records (including
292,666 benign packets), connected timing vectors, three X310 bandwidth runs,
four Data ablation rows, detector features/scores, and NLOS and antenna records.
These are experiment-specific workloads. The advertising half of Table~4 and
independent real-boundary PIP score reconstruction are outside this appendix's
reproduction scope.

\subsection{Set Up}
\subsubsection{Installation}
Extract both ZIPs into separate directories. In the source archive root run:
\begin{verbatim}
python3 -m venv .venv
.venv/bin/python -m pip install \
  -r requirements-test.txt
\end{verbatim}
In the experiment archive root run:
\begin{verbatim}
python3 -m venv .venv
.venv/bin/python -m pip install \
  -r requirements.txt
\end{verbatim}
Dependency installation needs network access or a local wheel cache;
subsequent analysis is offline. Commands below run from the corresponding
archive root, using its own virtual environment.

\subsubsection{Basic test}
In the source root, execute
\texttt{.venv/bin/python -m pytest -q tests}; the archived suite's expected
result is \texttt{163 passed}. In the experiment root, execute
\texttt{.venv/bin/python scripts/recompute\_tables.py}; successful completion
creates three CSVs in \file{outputs/recomputed_tables/}. A test pass checks
software functionality; claim validation uses the data and definitions below.

\subsection{Evaluation Workflow}
\subsubsection{Major claims}
\begin{compactitem}
\item[(C1)] The recorded evaluation supports receiver behavior across bandwidth,
range, SNR, and BPL/CAF configurations (paper evaluation, Fig.~10 and
Tables~3--4; E1). The antenna comparison measures recovered candidate frames
under a specified counting rule.
\item[(C2)] The retained timing samples support approximately unchanged
protocol-visible timing under the stated sample selection (paper
observability evaluation, Fig.~12 and Table~5; E2).
\item[(C3)] Retained IQ windows quantify spectral similarity; packet features
and session-held-out scores quantify tail-energy observability for the
supported detector views (paper Fig.~11 and countermeasure evaluation; E3).
\end{compactitem}

\subsubsection{Experiments}
For the complete offline workflow, run
\texttt{.venv/bin/python reproduce\_all.py} in the experiment root. Expect
15 \texttt{[OK]} lines and exit status 0, including standalone, regression,
plot, and metadata checks. A successful regression establishes agreement with
the archived baseline. Use the following interpretation of the results.
Time budgets below are planning estimates on a CPU workstation, excluding
downloads; allow 10 person-minutes for initial setup and up to 30 compute-minutes
for the complete run.

\begin{compactitem}
\item[(E1)] \textbf{Receiver and auxiliary replay (C1).}
\emph{Budget: 5 person-minutes + 10 compute-minutes.}
\emph{Preparation:} use the bundled range, bandwidth, ablation, and auxiliary
records without changing thresholds or excluding additional runs.
\emph{Execution:} run the complete workflow above.
\emph{Results:} inspect \file{figure10a.pdf}--\file{figure10c.pdf},
\file{outputs/recomputed_tables/table3_bandwidth.csv},
\file{table4_ablation.csv} in the same directory, and
\file{outputs/auxiliary_experiments/}. Table~3 aggregates supplied run-level
metrics; Table~4 covers four Data conditions. Range analysis applies sequence
deduplication, pattern-based BER/PSR, and its documented collision-outlier rule
($\mathrm{BER}>0.10$).

For the 60-m antenna replay, candidate recovery is the number of unique
sequence representatives divided by 200 nominal opportunities per run.
Full-frame rate counts the six-byte header/check and does not require a valid
XOR check. Across 11 default and 8 SAS-571 runs, the corresponding means are
55.0455\% / 53.9371\,kbps and 88.5625\% / 86.7857\,kbps.
These are candidate-recovery and full-frame-rate metrics, not verified useful
payload delivery. The companion \file{verify_antenna_correctness.py}, run
with the extracted experiment root as its argument after the replay,
applies exact-pattern plus integrity checking to the same selected
candidates, yielding 39.0909\% and 61.4375\%; excluding header/check and the fixed
A5 marker yields 37.2122 and 58.4856\,kbps. Nominal opportunities are not
independent transmitter ground truth. This correctness comparison is specific
to the antenna records and is not extrapolated to other experiments.

\item[(E2)] \textbf{Timing replay (C2).}
\emph{Budget: 5 person-minutes + 10 compute-minutes.}
\emph{Preparation:} retain the supplied packet groups and timing vectors.
\emph{Execution:} run the complete workflow; tables alone can be regenerated
with \texttt{.venv/bin/python scripts/recompute\_tables.py}.
\emph{Results:} inspect \file{figure12a.png}, \file{figure12b.png}, and
\file{outputs/recomputed_tables/table5_conn_timing.csv}.
Peripheral event-interval means are 29.999519, 29.999529, and 29.999582\,ms
for input datasets \texttt{normal}, \texttt{embed\_16B}, and
\texttt{embed\_240B}, with $n=59,55,57$; each median is 29.999971\,ms.
The central normal gap/event means are 150.300\,$\mu$s / 50.000906\,ms.
The upstream peripheral analyzer uses adjacent channel-block first-packet
times, a preceding block ending P$\to$C and containing an HR notification, and a
30$\pm$5\,ms interval filter. Conclusions concern these selected samples.
Use \texttt{Input\_dataset} to identify the short-tail condition:
the script's \texttt{embed\_8B} output alias maps to input
\texttt{embed\_16B}; the alias alone does not establish transmitted length.

\item[(E3)] \textbf{Spectral and detector replay (C3).}
\emph{Budget: 5 person-minutes + 10 compute-minutes.}
\emph{Preparation:} use the supplied IQ windows, feature eligibility masks,
and session-level out-of-fold scores.
\emph{Execution:} run the complete workflow; the detector-only entry is
\texttt{.venv/bin/python scripts/tail\_energy.py}.
\emph{Results:} quantitative spectral evidence is in
\file{outputs/frequency_similarity.csv} and
\file{outputs/frequency_pip_pairwise.csv}. Cosine similarity is computed on
Welch PSDs converted to dB and independently min--max normalized, at 4\,MS/s
with a 1024-point FFT. Adv uses windows 0022/0049, Data 0002/0000; PIP is
the mean of 24 pairs between normal windows 0000--0005 and PIP windows
0006--0009. The respective similarities are 99.7693\%, 98.8913\%, and
96.0272\%. A single spectrum overlay is not the 24-pair PIP statistic;
validation uses these numeric outputs.

Inspect \file{outputs/tail_energy/advertising_formal_roc.csv} and
\file{connected_session_oof.csv} in the same directory. Advertising AUC/TPR
is recomputed from eligible packet features. Connected five-fold session-level
scores cover direct-tail and PIP outer-boundary views; AUC / TPR at
FPR $\leq$5\% is 0.962703 / 75.6757\% and 0.520000 / 3.8462\%, respectively.
The separate \file{connected_pip_auc.csv} is an archived summary, not a
new score-level reconstruction. Its real-boundary 4/64\,$\mu$s guard/window
diagnostic (AUC 0.935128, TPR 71.1538\%) and the historical fixed 16/6\,$\mu$s
diagnostic (0.957949, 76.9231\%) have different parameters. Neither is claimed
here as independently reproduced session-held-out real-boundary evidence.
\end{compactitem}

\subsection{Notes on Reusability}
Use \file{DEPENDENCY_MAP.md} to replace inputs while preserving schemas,
grouping, sample masks, and metric definitions. Keep original records and
environment manifests with new outputs. The firmware guide provides optional
flash/control commands and expected acknowledgements; hardware/RF acceptance
requires a separate live run.

\subsection{Version}
Based on the LaTeX template for Artifact Evaluation V20220926.
\end{document}
